% Pitch final (10 a 15 minutos). Números e figuras de resultados entram sozinhos após os treinos.
\documentclass[aspectratio=169,10pt,t]{beamer}
\newcommand{\tipodoc}{Pitch}
% Tema comum dos slides (acompanhamento e pitch).
\usepackage{fontspec}
\setsansfont{Noto Sans}
\setmonofont[Scale=0.85]{Hack}
\usepackage[none]{hyphenat}
\usepackage{booktabs,tabularx,multirow,array,graphicx,xcolor}
\usepackage{tikz}
\usepackage{newunicodechar}
% Noto Sans não tem estes símbolos; usa a versão matemática.
\newunicodechar{→}{\ensuremath{\rightarrow}}
\newunicodechar{≈}{\ensuremath{\approx}}
\newunicodechar{≥}{\ensuremath{\geq}}
\newunicodechar{×}{\ensuremath{\times}}
\newunicodechar{↔}{\ensuremath{\leftrightarrow}}
\newcolumntype{L}{>{\raggedright\arraybackslash}X}
% Macros comuns ao artigo e aos slides.
\definecolor{azul}{HTML}{1B3A5C}
\definecolor{azulclaro}{HTML}{2E6B8F}
\definecolor{cinza}{HTML}{5F6B76}
\definecolor{pendcor}{HTML}{B3261E}
\definecolor{pendfundo}{HTML}{FBEDEC}

% Valores gerados por scripts/dados_docs.py; o que ainda não existe aparece como campo em branco.
\newcommand{\resdef}[2]{\expandafter\def\csname res@#1\endcsname{#2}}
\newcommand{\vazio}{\textcolor{pendcor}{\rule[-1.5pt]{2.4em}{0.5pt}}}
\newcommand{\res}[1]{\ifcsname res@#1\endcsname\csname res@#1\endcsname\else\vazio\fi}
\InputIfFileExists{gerado/dados.tex}{}{}

% Trecho a completar depois dos treinos (aparece em vermelho para não passar despercebido).
\newcommand{\pend}[1]{\textcolor{pendcor}{[#1]}}

% \figuraarq{largura}{altura}{arquivo em resultados/figuras}{o que gera a figura}
\newcommand{\figdir}{../../resultados/figuras}
\newcommand{\figuraarq}[4]{%
  \IfFileExists{\figdir/#3}%
    {\includegraphics[width=#1,height=#2,keepaspectratio]{\figdir/#3}}%
    {\fcolorbox{pendcor}{pendfundo}{\parbox[c][#2][c]{\dimexpr#1-2\fboxsep-2\fboxrule}{\centering\small
      \textcolor{pendcor}{\textbf{Figura pendente}}\\[2pt]\texttt{resultados/figuras/}\\\texttt{\detokenize{#3}}\\[2pt]#4}}}}


\renewcommand{\figurename}{Figura}
\renewcommand{\tablename}{Tabela}

\setbeamersize{text margin left=9mm,text margin right=9mm}
\setbeamertemplate{navigation symbols}{}
\setbeamercolor{normal text}{fg=black!85}
\setbeamercolor{structure}{fg=azul}
\setbeamercolor{alerted text}{fg=azulclaro}
\setbeamerfont{frametitle}{size=\Large,series=\bfseries}
\setbeamertemplate{itemize item}{\color{azulclaro}\raisebox{1.5pt}{\rule{4pt}{4pt}}}
\setbeamertemplate{itemize subitem}{\color{cinza}--}
\setbeamertemplate{enumerate item}{\color{azulclaro}\bfseries\insertenumlabel.}
\setlength{\leftmargini}{1.3em}

% Quem apresenta cada slide (participação equilibrada).
\newcommand{\falante}{}
\newcommand{\fala}[1]{\renewcommand{\falante}{#1}}

\setbeamertemplate{frametitle}{%
  \vskip6mm
  \begin{beamercolorbox}[wd=\paperwidth,leftskip=9mm,rightskip=9mm]{frametitle}
    \usebeamerfont{frametitle}\color{azul}\insertframetitle
    \ifx\falante\empty\else\hfill{\normalsize\mdseries\color{cinza}\falante}\fi
  \end{beamercolorbox}
  \vskip-1mm{\color{azulclaro}\hspace*{9mm}\rule{18mm}{1.6pt}}\vskip2mm}

\setbeamertemplate{footline}{%
  \hbox to\paperwidth{\hspace*{9mm}\color{cinza}\tiny Projeto 3 \textperiodcentered{} Classificação de Cenas Aéreas \textperiodcentered{} \tipodoc\hfill\insertframenumber\hspace*{9mm}}\vskip3mm}

\setbeamertemplate{title page}{%
  \vfill
  {\color{azulclaro}\small\bfseries CIÊNCIA DE DADOS \textperiodcentered{} UNIDADE III \textperiodcentered{} UNIFOR\par}
  \vskip3mm
  {\color{azul}\fontsize{26}{30}\selectfont\bfseries\inserttitle\par}
  \vskip2mm
  {\color{azulclaro}\large\insertsubtitle\par}
  \vskip3mm{\color{azulclaro}\rule{25mm}{2pt}}\vskip5mm
  {\normalsize\insertauthor\par}
  \vskip2mm
  {\small\color{cinza}\insertdate\par}
  \vfill}

\newcommand{\repo}{\href{https://github.com/JoeyAlanS/redes-reurais-fotos-satelite}{github.com/JoeyAlanS/redes-reurais-fotos-satelite}}
% Caixa de destaque
\newcommand{\destaque}[1]{\colorbox{azul!7}{\parbox{\dimexpr\linewidth-2\fboxsep}{\raggedright\small #1}}}


\title{Classificação de Cenas Aéreas com Redes Neurais}
\subtitle{MLP, CNN e o efeito da resolução no NWPU-RESISC45}
\author{Paulo Henrico Rabelo (2312652) \quad Joey Alan (2320416)}
\date{Projeto 3 \textperiodcentered{} \repo}

\begin{document}

\begin{frame}[plain]
\titlepage
\end{frame}

\fala{Paulo}
\begin{frame}{O problema}
\begin{columns}[T]
\begin{column}{0.52\textwidth}
\begin{itemize}
\item Classificar uma cena aérea inteira: aeroporto, floresta, porto…
\item Base para mapeamento urbano, monitoramento ambiental e planejamento
\item \textbf{Pergunta:} como a arquitetura e a resolução da imagem afetam o acerto, e quando o ganho compensa o custo?
\end{itemize}
\vskip3mm
\destaque{MLP vê a imagem como uma lista de 49.152 números.\\CNN vê vizinhanças: bordas, texturas, formas.}
\end{column}
\begin{column}{0.44\textwidth}
\centering\figuraarq{\linewidth}{5.2cm}{eda_exemplos.png}{notebooks/01\_exploracao.ipynb}
\end{column}
\end{columns}
\end{frame}

\begin{frame}{Dados: 10 classes do NWPU-RESISC45}
\begin{columns}[T]
\begin{column}{0.46\textwidth}
\begin{itemize}
\item \res{total} imagens RGB 256\,×\,256, 700 por classe (Google Earth, 30 a 0,2 m/pixel)
\item Ambientes naturais (floresta, deserto, montanha, rio, praia) e construídos (aeroporto, porto, residencial, tanques, agrícola circular)
\item \res{problemas} arquivos com problema; nenhuma exclusão
\item \res{dup/exatas} duplicatas exatas e \res{dup/quase} quase-duplicatas encontradas
\end{itemize}
\end{column}
\begin{column}{0.5\textwidth}
\centering\figuraarq{\linewidth}{5.2cm}{eda_duplicatas.png}{notebooks/01\_exploracao.ipynb}
\end{column}
\end{columns}
\end{frame}

\begin{frame}{Preparação sem vazamento}
\begin{columns}[T]
\begin{column}{0.48\textwidth}
\begin{enumerate}
\item SHA-256 + dHash (8 orientações) agrupam versões da mesma cena
\item Sorteio 70/15/15 por classe, semente 42, com cada grupo como unidade: \res{part/total/treino} / \res{part/total/validacao} / \res{part/total/teste}
\item Uma única divisão salva e reutilizada nos 6 modelos e no bônus
\item Pixels ÷ 255: nenhuma estatística aprendida
\item Aumento (espelhos e giros de 90°) só no treino
\end{enumerate}
\end{column}
\begin{column}{0.48\textwidth}
\destaque{\textbf{Teste intocado:} as escolhas usam só a validação. O teste foi avaliado uma única vez, no fim.}
\vskip3mm
\centering\figuraarq{\linewidth}{3.4cm}{eda_contagens.png}{notebooks/01\_exploracao.ipynb}
\end{column}
\end{columns}
\end{frame}

\begin{frame}{O que a exploração antecipou}
\begin{columns}[T]
\begin{column}{0.44\textwidth}
Em cor e textura global (só treino):
\begin{itemize}
\item Floresta e deserto: fáceis (89\% e 76\%)
\item Montanha (14\%): se mistura com rio, agrícola circular e aeroporto
\item Aeroporto: confunde com rio (22\%)
\end{itemize}
\vskip2mm
\destaque{Hipótese: forma e organização espacial são decisivas, e a CNN deve levar vantagem.}
\end{column}
\begin{column}{0.53\textwidth}
\centering\figuraarq{\linewidth}{5.2cm}{eda_ambiguidades.png}{notebooks/01\_exploracao.ipynb}
\end{column}
\end{columns}
\end{frame}

\fala{Joey}
\begin{frame}{Quatro arquiteturas, seis treinos}
\small
\begin{tabularx}{\linewidth}{@{}lLrL@{}}
\toprule
Modelo & Estrutura & Parâmetros & Hipótese\\
\midrule
MLP A & 49.152 → Dense 256 → 10 & 12,59 M & referência sem viés espacial\\
MLP B & 49.152 → 256 → 128 → 64 → 10 & 12,63 M & isola a \textbf{profundidade}\\
CNN A & 3 blocos Conv-BN-ReLU-Pool (32/64/128) → GAP → 10 & 95 mil & padrões locais, 132× menos parâmetros\\
CNN B & 4 blocos (até 256 filtros) → GAP → Dense 128 → 10 & 424 mil & campo receptivo 22 → 46 px\\
\bottomrule
\end{tabularx}
\vskip3mm
\begin{columns}[T]
\begin{column}{0.48\textwidth}
\destaque{\textbf{Treinos:} MLP-A/B-128, CNN-A/B-128 e CNN-A/B-256, todos com pesos novos.}
\end{column}
\begin{column}{0.48\textwidth}
\destaque{\textbf{GAP:} a mesma CNN tem os mesmos parâmetros em 128 e 256 px. Só a resolução muda.}
\end{column}
\end{columns}
\end{frame}

\begin{frame}{Protocolo}
\begin{columns}[T]
\begin{column}{0.5\textwidth}
\begin{itemize}
\item Adam (0,001), entropia cruzada, lote 32, semente 42
\item Até 40 épocas; para após 8 sem melhora
\item Checkpoint = maior \textbf{F1 macro de validação}
\item Mesmas condições para cada CNN nas duas resoluções
\item Hardware: \pend{GPU utilizada}
\end{itemize}
\end{column}
\begin{column}{0.46\textwidth}
\destaque{\textbf{Lição de percurso:} com o \textit{momentum} padrão do BatchNorm (0,99), a validação colapsava (11\% contra 66\% no treino). Com 0,9, as estatísticas de inferência acompanham o treino.}
\end{column}
\end{columns}
\end{frame}

\fala{Paulo}
\begin{frame}{Resultados no teste: desempenho e custo}
\centering\small
\begin{tabular}{@{}lrrrrrr@{}}
\toprule
Modelo & F1 valid. & Acurácia teste & F1 macro teste & Parâm. & Min de treino & s/época\\
\midrule
MLP-A-128 & \res{MLP-A-128/f1_macro_validacao} & \res{MLP-A-128/acuracia_teste} & \res{MLP-A-128/f1_macro_teste} & 12,59 M & \res{MLP-A-128/tempo_total_min} & \res{MLP-A-128/tempo_medio_epoca_s}\\
MLP-B-128 & \res{MLP-B-128/f1_macro_validacao} & \res{MLP-B-128/acuracia_teste} & \res{MLP-B-128/f1_macro_teste} & 12,63 M & \res{MLP-B-128/tempo_total_min} & \res{MLP-B-128/tempo_medio_epoca_s}\\
CNN-A-128 & \res{CNN-A-128/f1_macro_validacao} & \res{CNN-A-128/acuracia_teste} & \res{CNN-A-128/f1_macro_teste} & 95 mil & \res{CNN-A-128/tempo_total_min} & \res{CNN-A-128/tempo_medio_epoca_s}\\
CNN-B-128 & \res{CNN-B-128/f1_macro_validacao} & \res{CNN-B-128/acuracia_teste} & \res{CNN-B-128/f1_macro_teste} & 424 mil & \res{CNN-B-128/tempo_total_min} & \res{CNN-B-128/tempo_medio_epoca_s}\\
CNN-A-256 & \res{CNN-A-256/f1_macro_validacao} & \res{CNN-A-256/acuracia_teste} & \res{CNN-A-256/f1_macro_teste} & 95 mil & \res{CNN-A-256/tempo_total_min} & \res{CNN-A-256/tempo_medio_epoca_s}\\
CNN-B-256 & \res{CNN-B-256/f1_macro_validacao} & \res{CNN-B-256/acuracia_teste} & \res{CNN-B-256/f1_macro_teste} & 424 mil & \res{CNN-B-256/tempo_total_min} & \res{CNN-B-256/tempo_medio_epoca_s}\\
\bottomrule
\end{tabular}
\vskip4mm
\raggedright
\destaque{\pend{Mensagem principal em uma frase: quem venceu, por quanto e a que custo.}}
\end{frame}

\begin{frame}{Curvas de aprendizagem}
\centering\figuraarq{0.92\linewidth}{4.9cm}{painel_curvas_perda.png}{scripts/avaliar.py}
\vskip2mm
\raggedright\small\pend{Subajuste ou sobreajuste? Em que época está cada checkpoint? A MLP diverge cedo entre treino e validação?}
\end{frame}

\begin{frame}{MLP × CNN em 128\,×\,128}
\begin{columns}[T]
\begin{column}{0.38\textwidth}
\begin{itemize}
\item MLP A: \res{MLP-A-128/f1_macro_teste} \quad MLP B: \res{MLP-B-128/f1_macro_teste}
\item CNN A: \res{CNN-A-128/f1_macro_teste} \quad CNN B: \res{CNN-B-128/f1_macro_teste}
\item \pend{A profundidade ajudou a MLP? A CNN, com 132× menos parâmetros, venceu?}
\end{itemize}
\end{column}
\begin{column}{0.58\textwidth}
\centering\figuraarq{\linewidth}{5.2cm}{f1_por_classe_teste.png}{scripts/avaliar.py}
\end{column}
\end{columns}
\end{frame}

\fala{Joey}
\begin{frame}{Resolução: 128 × 256 e o custo}
\begin{columns}[T]
\begin{column}{0.38\textwidth}
\begin{itemize}
\item CNN A: \res{CNN-A-128/f1_macro_teste} → \res{CNN-A-256/f1_macro_teste} (F1 teste)
\item CNN B: \res{CNN-B-128/f1_macro_teste} → \res{CNN-B-256/f1_macro_teste}
\item Tempo/época A: \res{CNN-A-128/tempo_medio_epoca_s} → \res{CNN-A-256/tempo_medio_epoca_s} s
\item Tempo/época B: \res{CNN-B-128/tempo_medio_epoca_s} → \res{CNN-B-256/tempo_medio_epoca_s} s
\item \pend{Compensa? A CNN B ganhou mais, como previa o campo receptivo?}
\end{itemize}
\end{column}
\begin{column}{0.58\textwidth}
\centering\figuraarq{\linewidth}{5.2cm}{desempenho_custo_teste.png}{scripts/avaliar.py}
\end{column}
\end{columns}
\end{frame}

\begin{frame}{Onde os modelos se confundem}
\begin{columns}[T]
\begin{column}{0.3\textwidth}
\small
\pend{Top 3 pares confundidos e a explicação visual:}
\begin{itemize}
\item \pend{montanha ↔ …}
\item \pend{aeroporto ↔ …}
\item \pend{praia ↔ …}
\end{itemize}
\end{column}
\begin{column}{0.67\textwidth}
\centering\figuraarq{\linewidth}{5.6cm}{painel_confusao_teste.png}{scripts/avaliar.py}
\end{column}
\end{columns}
\end{frame}

\begin{frame}{Erros com alta confiança (\res{sel/modelo_erros})}
\centering\figuraarq{\linewidth}{3.9cm}{erros_teste.png}{scripts/avaliar.py}
\vskip3mm
\raggedright\small Modelo de maior F1 macro de validação; seis erros de teste com maior confiança.\\
\pend{O que cada erro revela: rótulo ambíguo, contexto dominante ou detalhe perdido?}
\end{frame}

\fala{Paulo}
\begin{frame}{Limitações}
\begin{itemize}
\item \textbf{Um treino por configuração:} diferenças pequenas podem ser só variação entre sementes
\item \textbf{Mesma distribuição:} o teste vem do mesmo Google Earth, não de novas regiões, sensores ou estações
\item \textbf{Quase-duplicatas:} o \textit{hash} não pega cenas vizinhas parcialmente sobrepostas, então o teste pode ser otimista
\item \textbf{Resolução simulada:} reduzimos a mesma imagem; um sensor mais grosseiro real tem outro ruído e outro borramento
\item \textbf{Sem pré-treino} nos 6 modelos (exigência do projeto); o bônus YOLO mostra o efeito do pré-treino
\end{itemize}
\end{frame}

\fala{Paulo e Joey}
\begin{frame}{Conclusões}
\begin{enumerate}
\item \pend{MLP × CNN: …}
\item \pend{Resolução: …}
\item \pend{Custo-benefício: …}
\item \pend{Classes difíceis: …}
\end{enumerate}
\vskip4mm
\destaque{\textbf{Próximos passos:} repetir com várias sementes, testar pesos pré-treinados e validar em outra base (por exemplo, AID).}
\vfill
\centering\small Código, partições, métricas e figuras: \repo
\end{frame}

\begin{frame}[plain]
\vfill
\centering
{\color{azul}\fontsize{30}{34}\selectfont\bfseries Obrigado!\par}
\vskip4mm
{\large\color{azulclaro}Perguntas?\par}
\vskip6mm
\small Paulo Henrico Rabelo (2312652) \quad Joey Alan (2320416)\\[2pt]
\repo
\vfill
\end{frame}

\end{document}
